\begin{ex}
Một hình vuông có cạnh dài 8 cm. Chu vi của hình vuông đó là bao nhiêu xăng-ti-mét?
\shortans{32}
\loigiai{
$8 \times 4 = 32$ cm.
}
\end{ex}

\begin{ex}
Một hình vuông có cạnh dài 8 cm. Diện tích của hình vuông đó là bao nhiêu xăng-ti-mét vuông?
\shortans{64}
\loigiai{
$8 \times 8 = 64\text{ cm}^2$.
}
\end{ex}

\begin{ex}
Một hình chữ nhật có chiều dài 12 cm, chiều rộng 7 cm. Chu vi hình chữ nhật đó là bao nhiêu xăng-ti-mét?
\shortans{38}
\loigiai{
$(12 + 7) \times 2 = 38$ cm.
}
\end{ex}

\begin{ex}
Một hình chữ nhật có chiều dài 12 cm, chiều rộng 7 cm. Diện tích hình chữ nhật đó là bao nhiêu xăng-ti-mét vuông?
\shortans{84}
\loigiai{
$12 \times 7 = 84\text{ cm}^2$.
}
\end{ex}

\begin{ex}
Một hình tròn có bán kính 8 cm. Đường kính của hình tròn đó là bao nhiêu xăng-ti-mét?
\shortans{16}
\loigiai{
$8 \times 2 = 16$ cm.
}
\end{ex}

\begin{ex}
Một hình tròn có đường kính 20 cm. Bán kính của hình tròn đó là bao nhiêu xăng-ti-mét?
\shortans{10}
\loigiai{
$20 : 2 = 10$ cm.
}
\end{ex}

\begin{ex}
Khối lập phương có bao nhiêu mặt?
\shortans{6}
\loigiai{
Có 6 mặt.
}
\end{ex}

\begin{ex}
Khối lập phương có bao nhiêu đỉnh?
\shortans{8}
\loigiai{
Có 8 đỉnh.
}
\end{ex}

\begin{ex}
Khối hộp chữ nhật có bao nhiêu cạnh?
\shortans{12}
\loigiai{
Có 12 cạnh.
}
\end{ex}

\begin{ex}
Đổi: $5\text{ m } 8\text{ cm} = \dots\text{ cm}$.
\shortans{508}
\loigiai{
$500 + 8 = 508$ cm.
}
\end{ex}

\begin{ex}
Đổi: $3\text{ km} = \dots\text{ m}$.
\shortans{3000}
\loigiai{
$3\text{ km} = 3000\text{ m}$.
}
\end{ex}

\begin{ex}
Đổi: $1\text{ kg} = \dots\text{ g}$.
\shortans{1000}
\loigiai{
$1\text{ kg} = 1000\text{ g}$.
}
\end{ex}

\begin{ex}
Đổi: $1\text{ l} = \dots\text{ ml}$.
\shortans{1000}
\loigiai{
$1\text{ l} = 1000\text{ ml}$.
}
\end{ex}

\begin{ex}
Một hình vuông có chu vi 40 cm. Cạnh của hình vuông đó dài bao nhiêu xăng-ti-mét?
\shortans{10}
\loigiai{
$40 : 4 = 10$ cm.
}
\end{ex}

\begin{ex}
Một hình vuông có diện tích $49\text{ cm}^2$. Cạnh của hình vuông đó dài bao nhiêu xăng-ti-mét?
\shortans{7}
\loigiai{
Vì $7 \times 7 = 49$ nên cạnh là 7 cm.
}
\end{ex}

\begin{ex}
Một hình chữ nhật có diện tích $60\text{ cm}^2$, chiều rộng 5 cm. Chiều dài hình chữ nhật đó là bao nhiêu xăng-ti-mét?
\shortans{12}
\loigiai{
$60 : 5 = 12$ cm.
}
\end{ex}

\begin{ex}
Nhiệt độ sôi của nước ở điều kiện chuẩn là bao nhiêu độ C?
\shortans{100}
\loigiai{
$100^\circ$C.
}
\end{ex}

\begin{ex}
Nhiệt độ nước đá đang tan là bao nhiêu độ C?
\shortans{0}
\loigiai{
$0^\circ$C.
}
\end{ex}

\begin{ex}
Một gói bánh nặng 250 g, 4 gói bánh như thế nặng bao nhiêu gam?
\shortans{1000}
\loigiai{
$250 \times 4 = 1000$ g.
}
\end{ex}

\begin{ex}
Chai dầu có 750 ml, dùng hết 250 ml. Chai dầu còn lại bao nhiêu mi-li-lít?
\shortans{500}
\loigiai{
$750 - 250 = 500$ ml.
}
\end{ex}

\begin{ex}
Đoạn thẳng AB dài 18 cm, M là trung điểm của AB. Độ dài đoạn AM là bao nhiêu xăng-ti-mét?
\shortans{9}
\loigiai{
$18 : 2 = 9$ cm.
}
\end{ex}

\begin{ex}
Một góc vuông có số đo là bao nhiêu độ?
\shortans{90}
\loigiai{
Góc vuông bằng $90^\circ$.
}
\end{ex}

\begin{ex}
Chu vi hình tam giác đều cạnh 15 cm là bao nhiêu xăng-ti-mét?
\shortans{45}
\loigiai{
$15 \times 3 = 45$ cm.
}
\end{ex}

\begin{ex}
Chu vi hình tứ giác có 4 cạnh đều bằng 12 cm là bao nhiêu xăng-ti-mét?
\shortans{48}
\loigiai{
$12 \times 4 = 48$ cm.
}
\end{ex}

\begin{ex}
Một mảnh vườn hình chữ nhật có chiều dài 20 m, rộng 10 m. Diện tích mảnh vườn là bao nhiêu mét vuông?
\shortans{200}
\loigiai{
$20 \times 10 = 200\text{ m}^2$.
}
\end{ex}
